
\chapter{Analysis of the legacy system}

\section{Technologies used}
\begin{itemize}
\item description of microsoft access 
\item the advantages and diadvantages of using microsoft access as a DBMS
\item why bfg has to move to something else
\item the new requirements that need to be met
\end{itemize}

\section{The current database layout}
\begin{itemize}
\item the format I recieved the data in
\item some assumptions made, and why I had to make them
\end{itemize}

\subsection{Considerations for the new database}
\begin{itemize}
\item taxonomy, how the fungi lists need to be formated to account for this
\item the problem with the fungi dictionary being outdated, and why it can't be updated
\item lots of unnessassary fields are no longer needed however should be archived still
\item new requirements
\end{itemize}

\subsection{Bad database practices}
\begin{itemize}
\item certain fields e.g substrates had predefined lists that weren't always used when creating records, this led to lots of data redundancy
\item the format for the fungi dictionary wasn't always followed leading to data redundancy
\item format for the site grid refs weren't consistent, making it harder to perform calculations on them
\item normalisation problems, some fields had data that was the same as some other fields, e.g a fungus fullName field was just a concatination of the genus, species, and varity
\end{itemize}